% Nouveau EGA V (cote 25): archive cote 25, site page 145 (PDF page 146).
% typescript. Draft compared: none.
% AI-transcribed by Claude Opus 5.5 (claude-opus-5-5, adaptive thinking, effort high), Message Batches API, batch ega25b (msgbatch_01A7n5YuXubUPNK1zsepi8RG), 2026-11-13; not yet checked by eye.
\scanpage{25/145}
\margin{6}

\uline{Proposition} 7. Sous les conditions précédentes (à rappeler), l'ensemble
sous-jacent à $V(\phi)^{\text{sup sing}}$ n'est autre que l'ensemble des points de
$V(\phi)^{\text{sing}}$ \xout{\illegible{1}} où le morphisme $V(\phi)^{\text{sing}}\longrightarrow Y$ (de préschémas lisses
sur S de dimension relative m-1 et m respectivement) est \uline{ramffié}.

En termes de papa (qu'il faudrait signaler en remarque), un point
$x \in V(\phi)$ est donc supersingulier rel à Y sss ``il consiste en au
moins deux points singuliers confondus'' ....

On peut et doit préciser prop.7 du point de vue \inserted{d'une identité de} \xout{des} sous-présché-%
mas et non seulement de sous-ensembles. En effet, $V(\phi)^{\text{sup sing}}$ a été
défini comme un sous-préschéma fermé de X , or on peut également définir
un sous-préschéma fermé naturel de $V^1$ de telle façon que l'ensemble
sous-jacent soit l'ensemble des points de ramification par rapport à Y.
En effet, il faut exprimer l'ensemble des points où un certain homomor-%
phisme de modules loc libres $Q$ ( $=\underline{\Omega}^1_{Y/S}\otimes\underline{O}_{V^1}$ ) $\longrightarrow$ $R$ ( $=\underline{\Omega}^1_{V^1/S}$ )
n'est pas surjectif. Si q et \unsure{m|r} sont les rang respectifs, c'est aussi
l'ensemble d es points où $\bigwedge^{\unsure{m|r}} Q \rightarrow \bigwedge^{\unsure{m|r}} R$ n'est pas surjectif, ce qui est
aussi l'ensemble des zéros de la section évidente de \uline{Hom} ( $\bigwedge^{\unsure{m|r}} Q$ , $\bigwedge^{\unsure{m|r}} R$)
$\simeq (\bigwedge^{q-r} Q)\otimes(\bigwedge^{r} R)$\hand{\illegible{1}}$\otimes(\bigwedge^{q} Q)$\hand{$^{\vee}$} , donc l'ensemble sous-jacent au \inserted{sous-}préschéma
fermé des zéros de cette section, soit \uline{Ram}($V^1/Y$) . Je dis que ce dernier
sous-schéma est identique à $V(\phi)^{\text{sup sing}}$ . C'est un simple exercice
sur le diagramme plus haut, compte tenu que $V(\phi)^{\text{sup sing}}$ est défini
par le même procédé que celui explicité pour $Q\longrightarrow R$, mais en termes de
l'homomorphisme $P$ ( $= \overset{\vee}{\underline{P}}{}^1_{X/Y}\otimes\underline{O}_{V^1}$ ) $\longrightarrow$ $S$ ($=\underline{\Omega}^1_{X/Y}\otimes\underline{O}_{V^1}$) , comme
il résulte de la description de \hand{$\mu$} donnée plus haut. On est alors
ramené à la situation générale suivante:

On a sur un \deleted{préschéma} \inserted{\unsure{espace annelé}} W un module localement libre \inserted{M} de rang m,
et deux sous-modules localement libres P et Q de rangs respectifs

\mspage{145}
